% GENERATED FILE. Edit canonical lessons, then run npm run generate:books.
% canonical-source-hash: fnv1a64:fcde80d6cc8d2932
% canonical-lessons: ZH-C420-jiazhao, ZH-C420-chedao, ZH-C420-chewei, ZH-C420-daoche, ZH-C420-weizhang

\chapter{Parking}
\label{ch:zh-parking}
\hlchaptermodality{420}

\begin{chapteropening}
\textbf{By the end of this chapter:} \emph{I can talk about a driving licence, traffic lanes, parking spaces, reversing and breaking the rules.}

Complete the last lesson of chapter 420: I can talk about a driving licence, traffic lanes, parking spaces, reversing and breaking the rules.
\end{chapteropening}

\section[jiàzhào]{jiàzhào — driving licence}

\label{lesson:ZH-C420-jiazhao}



\begin{quote}
Before the new one: say the Mandarin for international, then the Mandarin for inside the country.
\end{quote}

\subsection*{You'll want to know: jiàzhào}
\emph{jiàzhào} — \textquotedblleft{}driving licence\textquotedblright{}.

\begin{quote}
\emph{wǒ de jiàzhào bú jiàn le} — My driving licence has gone missing.
\end{quote}

Its characters come later; for now it is a word for the ear.

\begin{cousinweb}
\emph{jiàzhào} is built from \emph{jià} 'to drive' and \emph{zhào}, short for a word for a licence; both are inherited from Middle Chinese, \emph{kaeH} and \emph{tsyewH} (Guangyun, in Baxter's transcription), which glosses the first 'to drive a team' and the second 'bright'.
\end{cousinweb}

\begin{grammarlens}[title={What you've built}]
One more word for getting around.
\end{grammarlens}

\subsection*{Your turn}
\begin{itemize}
\raggedright
  \item \emph{Say it:} \emph{jiàzhào}
  \item \emph{Say it:} \emph{jiàzhào}, once more
  \item \emph{Say it:} \emph{wǒ de jiàzhào bú jiàn le} — My driving licence has gone missing.
  \item \emph{Recall:} say the Mandarin for international, then the Mandarin for inside the country, then say \emph{jiàzhào} again
\end{itemize}

\begin{tcolorbox}[breakable,colback=teal!4,colframe=teal!35!black,title={Before you move on}]
What does jiàzhào mean? (\textquotedblleft{}Driving licence\textquotedblright{}.) Say it once more.
\end{tcolorbox}
\section[\texorpdfstring{\zh{车道} (chēdào)}{车道 (chēdào)}]{\zh{车道} (chēdào) — traffic lane}

\label{lesson:ZH-C420-chedao}



\begin{quote}
Before the new one: say the Mandarin for inside the country, then the Mandarin for a driving licence.
\end{quote}

\subsection*{You'll want to know: \zh{车道}}
\textbf{\zh{车道}} — \emph{chēdào} — \textquotedblleft{}traffic lane\textquotedblright{}.

\begin{quote}
\emph{wǎng yòubian de chēdào kāi} — Move over into the right-hand lane.
\end{quote}

\begin{cousinweb}
\textbf{\zh{车道}} (\emph{chēdào}) is built from \zh{车} 'vehicle' and \zh{道} 'a road', both inherited from Middle Chinese, \emph{tsyhae} and \emph{dawX} (Guangyun, in Baxter's transcription), which glosses the second 'a road' and 'reason'.
\end{cousinweb}

\begin{grammarlens}[title={What you've built}]
One more word for getting around.
\end{grammarlens}

\subsection*{Your turn}
\begin{itemize}
\raggedright
  \item \emph{Say it:} \emph{chēdào}
  \item \emph{Say it:} \emph{chēdào}, once more
  \item \emph{Say it:} \emph{wǎng yòubian de chēdào kāi} — Move over into the right-hand lane.
  \item \emph{Recall:} say the Mandarin for inside the country, then the Mandarin for a driving licence, then say \emph{chēdào} again
\end{itemize}

\begin{tcolorbox}[breakable,colback=teal!4,colframe=teal!35!black,title={Before you move on}]
What does \zh{车道} mean? (\textquotedblleft{}Traffic lane\textquotedblright{}.) Say it once more.
\end{tcolorbox}
\section[chēwèi]{chēwèi — parking space}

\label{lesson:ZH-C420-chewei}



\begin{quote}
Before the new one: say the Mandarin for a driving licence, then the Mandarin for a traffic lane.
\end{quote}

\subsection*{You'll want to know: chēwèi}
\emph{chēwèi} — \textquotedblleft{}parking space\textquotedblright{}.

\begin{quote}
\emph{wǒ zài tíngchēchǎng zhǎo bu dào chēwèi} — I can't find a space in the car park.
\end{quote}

Its characters come later; for now it is a word for the ear.

\begin{cousinweb}
\emph{chēwèi} is built from \zh{车} 'vehicle' and \emph{wèi} 'a place', both inherited from Middle Chinese, \emph{tsyhae} and \emph{hwijH} (Guangyun, in Baxter's transcription).
\end{cousinweb}

\begin{grammarlens}[title={What you've built}]
One more word for getting around.
\end{grammarlens}

\subsection*{Your turn}
\begin{itemize}
\raggedright
  \item \emph{Say it:} \emph{chēwèi}
  \item \emph{Say it:} \emph{chēwèi}, once more
  \item \emph{Say it:} \emph{wǒ zài tíngchēchǎng zhǎo bu dào chēwèi} — I can't find a space in the car park.
  \item \emph{Recall:} say the Mandarin for a driving licence, then the Mandarin for a traffic lane, then say \emph{chēwèi} again
\end{itemize}

\begin{tcolorbox}[breakable,colback=teal!4,colframe=teal!35!black,title={Before you move on}]
What does chēwèi mean? (\textquotedblleft{}Parking space\textquotedblright{}.) Say it once more.
\end{tcolorbox}
\section[dàochē]{dàochē — to reverse the car; to back up}

\label{lesson:ZH-C420-daoche}



\begin{quote}
Before the new one: say the Mandarin for a traffic lane, then the Mandarin for a parking space.
\end{quote}

\subsection*{You'll want to know: dàochē}
\emph{dàochē} — \textquotedblleft{}to reverse the car; to back up\textquotedblright{}.

\begin{quote}
\emph{tā zài tíngchēchǎng dàochē} — He's reversing in the car park.
\end{quote}

Its characters come later; for now it is a word for the ear.

\begin{cousinweb}
\emph{dàochē} is built from \emph{dào} 'backwards' and \zh{车} 'vehicle', both inherited from Middle Chinese, \emph{tawH} and \emph{tsyhae} (Guangyun, in Baxter's transcription), which glosses the first as hanging upside down.
\end{cousinweb}

\begin{grammarlens}[title={What you've built}]
One more word for getting around.
\end{grammarlens}

\subsection*{Your turn}
\begin{itemize}
\raggedright
  \item \emph{Say it:} \emph{dàochē}
  \item \emph{Say it:} \emph{dàochē}, once more
  \item \emph{Say it:} \emph{tā zài tíngchēchǎng dàochē} — He's reversing in the car park.
  \item \emph{Recall:} say the Mandarin for a traffic lane, then the Mandarin for a parking space, then say \emph{dàochē} again
\end{itemize}

\begin{tcolorbox}[breakable,colback=teal!4,colframe=teal!35!black,title={Before you move on}]
What does dàochē mean? (\textquotedblleft{}To reverse the car; to back up\textquotedblright{}.) Say it once more.
\end{tcolorbox}
\section[wéizhāng]{wéizhāng — to break the traffic rules}

\label{lesson:ZH-C420-weizhang}



\begin{quote}
Before the new one: say the Mandarin for a parking space, then the Mandarin for to reverse the car.
\end{quote}

\subsection*{You'll want to know: wéizhāng}
\emph{wéizhāng} — \textquotedblleft{}to break the traffic rules\textquotedblright{}.

\begin{quote}
\emph{zhèr bù néng tíngchē tíng le jiù wéizhāng} — You can't park here: if you do, you're breaking the rules.
\end{quote}

Its characters come later; for now it is a word for the ear.

\begin{cousinweb}
\emph{wéizhāng} is built from \emph{wéi} 'to go against' and \emph{zhāng} 'a rule', both inherited from Middle Chinese, \emph{hjw+j} and \emph{tsyang} (Guangyun, in Baxter's transcription), which glosses the first 'to turn your back on' and the second as a section of writing; the sense 'a rule' is not given there.
\end{cousinweb}

\begin{grammarlens}[title={What you've built}]
One more word for getting around.
\end{grammarlens}

\subsection*{Your turn}
\begin{itemize}
\raggedright
  \item \emph{Say it:} \emph{wéizhāng}
  \item \emph{Say it:} \emph{wéizhāng}, once more
  \item \emph{Say it:} \emph{zhèr bù néng tíngchē tíng le jiù wéizhāng} — You can't park here: if you do, you're breaking the rules.
  \item \emph{Recall:} say the Mandarin for a parking space, then the Mandarin for to reverse the car, then say \emph{wéizhāng} again
\end{itemize}

\begin{tcolorbox}[breakable,colback=teal!4,colframe=teal!35!black,title={Before you move on}]
What does wéizhāng mean? (\textquotedblleft{}To break the traffic rules\textquotedblright{}.) Say it once more.
\end{tcolorbox}
